\chapter{Theoretical and methodological approaches}


This thesis has the objective to explore the use of amulets in Napatan Nubia through a multi-relational framework, taking into consideration variations in social, chronological and geographical contexts. To achieve this, I will employ theoretical and methodological models to interpret the data systematically.


In this chapter, I will begin by examining the initial interpretations of Nubian material culture, guided by the concept of "Egyptianization". This perspective emphasizes the influence of Egyptian culture on Nubians, viewed as a mere imitation. However, shifts in scholarship have led to the introduction of the idea of cultural entanglements, which highlights the dynamic interaction between Egyptian and Nubian cultures over millennia.


This shift underscores how Egyptian cultural features were adopted, rejected and adapted to meet local needs by the Nubians. To further enhance our understanding of amulet use during the Napatan period, I will introduce the concept of objectscapes, which has its origins in the criticism of Romanization. This framework provided an appropriate lens for interpreting the role of amulets within a particular context (social, regional and temporal).


Additionally, I will address methodological considerations regarding the chronological division of the Napatan period used in this thesis and the categories employed to classify the other objects found in the tombs. Since there were numerous types of objects, they were grouped into wider categories for better visualization and interpretation. Furthermore, I will outline the construction of the database and discuss the limitations and potentials of the available data.
This chapter will, thus, provide a solid foundation for the analysis and interpretations of the data in the following the chapters.


\section{From Egyptianization to cultural entanglement}


The concept of "Egyptianization" was first introduced by \textcite{reisnerArchaeologicalSurveyNubia1910}, when describing the similarities between Nubian and Egyptian graves from the New Kingdom. He specifically attributed this phenomenon to the immigration of Egyptians into Nubia due to the Hyksos rule during the Second Intermediate Period. Later, Reisner employed the term to describe the acculturation of Nubians into the Egyptian culture, positing that the decrease in indigenous material culture was the result of the inferior Nubian population passively adopting the superior Egyptian culture \parencites[p.~236-237]{reisnerOutlineHistorySudan1918}{vanpeltRevisingEgyptoNubianRelations2013}\parencite{weglarzContinuityChangeReevaluation2017}. For example, \textcite{bietakStudienZurChronologie1968} deduced that the C-Group submitted without resistance to Egyptian culture as the local burial customs were replaced by Egyptian ones, while \textcite[p.~10]{säve-söderberghScandinavianJointExpedition1989} described C-Group building techniques as "imitating Egyptian fortifications". These ideas were in accordance with theories of cultural diffusionism and acculturation, popular at the time, meaning that the migration of Egyptians into the Nubian territory would substitute the local culture in a monodirectional way \parencites{raueIntroduction2019}. This idea was especially applied to the New Kingdom conquest of Lower Nubia, where many aspects of local life, such as administrative styles, material culture and customs, were replaced by Egyptian forms \parencites[p.~523]{vanpeltRevisingEgyptoNubianRelations2013}. This interpretation was reinforced by the ancient Egyptians themselves, whose records shows that they considered themselves superior to foreigners. Different terminology in the ancient Egyptian language were used for foreigners and Nubians were described as "wretched" \parencites[p.~12-13]{smithWretchedKushEthnic2003}.


Elite emulation was applied in later models of Egyptianization. It was first used by Higginbotham in the context of the Egyptian occupation of the Levant and later applied to Nubia \parencites{higginbothamEgyptianizationEliteEmulation2000}. In short, it argues that the Egyptian administration co-opted elites into acculturation for political and social gain, which, in turn, supported the Egyptian rule of the region \parencites[p.~277-278]{smithNubiaEgyptInteraction1998}[p.~274]{törökTwoWorldsFrontier2009}. This idea is intricately linked with the core-periphery model, where the core controls the periphery, whose role is primarily as supply of resources. In this system, Nubia would provide Egypt with its natural resources, such as gold and ivory, while it relies on the dominant core for social, cultural and political influence \parencites{smithNubiaEgyptInteraction1998, flamminiAncientCorePeripheryInteractions2008}. This idea is not without merit, as during the New Kingdom rule in Lower Nubia, foreign children were educated in Egypt in order to become local rulers subject to the Egyptian administration \parencites[p.~530]{vanpeltRevisingEgyptoNubianRelations2013}[p.~223]{morrisAncientEgyptianImperialism2018}. This model, however, fails to consider the whole of Nubian society and how the lower levels of society were not the main subjects of Egyptian acculturation, since they had nothing to gain from it. Therefore, it was necessary to consider localized and bottom-up perspectives in order to understand the presence of Egyptian-style objects in non-elite contexts \parencites[p.~530]{vanpeltRevisingEgyptoNubianRelations2013}.


Postcolonial theories have been increasingly applied in archaeology to study the complex dynamic of ancient societies under colonial rule and highlighting local agency \parencites{gosdenPostcolonialArchaeologyIssues2001, vandommelenColonialMattersMaterial2006}. Following the critiques of the concept of Romanization influenced by postcolonial theory, \textcites{vanpeltRevisingEgyptoNubianRelations2013} proposed an approach, following Bhabha's idea of hybridity and "third space" \textcite{bhabhaLocationCulture1994} as re-assessed by \textcites{stockhammerConceptualizingCulturalHybridization2012, stockhammerHybridityEntanglementEssentialism2013}, centered around the concept of "cultural entanglements". According to the author, the concept aims to "identify and highlight archaeologically the creative potential of 'liminal spaces' which are conceptualized as situations and spaces of intercultural encounter", which can be achieved by emphasizing local agency and cultural practices \parencites[p.~533]{vanpeltRevisingEgyptoNubianRelations2013}. Cultural entanglements aims to provide a more bottom-up approach and consider "cultural change in a multidirectional and differentially negotiated in different moments and spaces of encounter" and, although, not denying Egyptian influence, does not privilege one culture over another \parencites[p.~541]{vanpeltRevisingEgyptoNubianRelations2013}. Following van Dommelen's cultural entanglements approach and Dietler's consumption model, Smith \textcite[p.~372]{smithNubianExperienceEgyptian2020} argues that individual choices shaped cultural practices, actively constructing a new configuration based on adoption, rejection and adaptation \parencites{vandommelenColonialInteractionsHybrid2005, dietlerConsumption2010}. The author further proposes that "through a nuanced analysis of the intersection of the different social logic of the parties involved, agency is both potentially discernible and historically crucial to understanding larger developments like the spread of Egyptian funerary practices in both Lower and Upper Nubia during the New Kingdom." 


Recently, this approach has been used in analyses of excavations in Nubia sites and cemeteries. However, the majority of these studies have focused on New Kingdom Nubia, where the dynamic between Egypt and Kush was of a colonial nature. Moreover, they fail to address the inherent violence of the colonizer in this context \parencite[p.~72-73]{matićPostcolonialismReverseDiscourse2023}. As a matter of fact, Matić posits that the concepts of hybridization, third space and cultural entanglements have been misused in Egyptian and Nubian archaeology due to misunderstanding of these ideas \parencite[p.~61, 74]{matićPostcolonialismReverseDiscourse2023}. In order to better make use of these concepts, the author advocates for "collection and study of data triggered by novel research questions" \parencite[p.~74]{matićPostcolonialismReverseDiscourse2023}. Nevertheless, the cultural entanglements approach has also failed to take into account the internal diversity within Nubia \parencites[p.~5]{lemosMaterialCultureColonization2020}. It is important to consider local consumption of objects through the diverse interests and practices of different social groups \parencites[p.~57]{dietlerConsumption2010}. The Napatan period features diverse social dynamics as well as regional and temporal variations that should be considered inside their own context. Furthermore, cultural interactions with Egypt shifted considerably throughout the century, with certain phases, such as the 25th Dynasty and Early Napatan, where the Kushites were in control of Upper Egypt and its institutions. Thus, use of amulets, has to be analyzed in each context, in order to illuminate local choices and agency. Therefore, it is important to consider other theoretical frameworks that contemplate this variability.


\section{Objectscapes}


Recent critiques of Roman imperialism and traditional "Romanization" approaches have introduced the concept of objectscapes, defined as "the material and stylistic properties of a repertoire of objects in a given period and geographic range", essentially a chart of a specific time and space \parencites[p.~368]{pittsObjectscapesManifestoInvestigating2021}. This approach emphasizes the circulation of objects within and across societies, highlighting how they actively shaped cultural practices. As such, it is particularly suitable for analyzing contexts exposed to trade and cultural exchanges \parencites[p.~7]{pittsRomanObjectRevolution2019}. Through their material and stylistics features, objects within an objectscape reshape human behaviour, though their effects vary between contexts \parencites[p.~197]{versluysObjectscapesMaterialConstitution2017}. Furthermore, the concept directly challenges simplistic notions of cultural homogenization, stressing the crucial role of adaptation and transformation as foreign goods integrate local contexts in order to fulfill specific social roles \parencites[p.~13-14]{pittsRomanObjectRevolution2019}. 


At the center of objectscape analysis lie the relational and cultural dynamics of material culture, examined through objects' specific contexts and interconnections. As Pitts and Versluys argue, "privileging internationality fosters better understanding of what objects \textit{did} in the past, and confronts the assumptions inherent in many studies in which objects are reduced to proxies for abstract processes (e.g. economic growth) or social categories (e.g. ethnicities and identities)" \parencite[p.~368]{pittsObjectscapesManifestoInvestigating2021}. In order to implement this approach, the authors identify three interconnected principles: connectivity, relationality, and impact. Connectivity concerns the origin and introduction of objects---where they came from and which ones were new to a given setting. Relationality addressees how objects relate to one another within a shared repertoire. Impact, meanwhile, captures the broader consequences of material change: how objects innovate, how they affect people and other objects, and how these effects unfold over both the short and long term \parencite[p.~369-371]{pittsObjectscapesManifestoInvestigating2021}. Together, these three principles reveal objectscapes not as static, but rather as active forces shaping historical change.


Pitts elaborates this framework with three additional concepts that build upon the previous ones: stylistic genealogy, local agency and replication, and longer-term evolution (tracking changes across time). The first concept focuses on the tracing of design influences in mass produced (not handmade) objects used for many years in a certain society. Thus, stylistic genealogy has the potential to indicate the way in which objects retained certain traits and embodied innovations at the same time, as well as their choice of design and function in a certain context. Local agency and replication examines how foreign mobile objects achieved significant cultural impact within local contexts, leading to local reproductions. This process also explores the use of objects and their distribution across various local settings. Object reproductions not only reflect adaptation of foreign elements but also reveal local production processes and the agency of artisans. Finally, longer-term evolution explores the impact of objects throughout many generations and not only in a particular moment. This evolution can be observed by tracing changes in material designs of a particular object across time, underscoring local agency in the shaping of objects \parencites[p.~16-19]{pittsRomanObjectRevolution2019}.


Together, these concepts provide archaeologists with a methodological road map for understanding foreign material culture in local contexts. They allow archaeological research to move beyond descriptions towards a more holistic interpretation of cultural exchange. By linking the historical roots of an object to its social recontextualization and its long-term developmental trajectory, the objectscapes framework avoids explanations rooted in diffusionism and acculturation (such as Egyptianization), where changes are understood as imposed from the outside. Moreover, it acknowledges that objects are simultaneously shaped by inherited traditions, the agency of local artisans, and the changing societal needs over time. This integrative approach reframes material culture as a dynamic participant in the ongoing negotiation of cultural identity. In contrast to the cultural entanglement approaches---which often rely on concepts of hybridity and "third space" that, as Matić warns are frequently misunderstood---objectscapes offers a more methodologically and critically grounded alternative based on material criteria that can be observed and through data. Therefore, objects in high connectivity contexts can be analyzed through Pitts and Versluys' key concepts in order to reconstruct objectscapes across different localities and periods. In this context, Bhabha's theory of hybridization proves particularly useful as well. Discussing Bhabha's work, \textcite{fahlanderThirdSpaceEncounters2007} emphasizes aspects applicable to archaeological research. For instance, the author addresses the idea that cultures are not homogeneous but composed of many different internal elements and contexts. Consequently, social encounters between distinct cultures are not limited to a simple blend of elements. In fact, they actually produce new and distinctive features that are incorporated over the long term.


Furthermore, the objectscapes framework aligns with recent discussions of globalization in the ancient world. Contrary to the popular idea that globalization can only describe consumption of foreign goods in the modern era, the term has been increasingly applied to the similar phenomena in the past \parencites{pittsGlobalisationCirculationMass2014}. Globalization is understood as "more than just another way to describe connectivity, trade, migration, internationalism, or diffusion", referring instead to active processes of increasing interconnection that were perceived by the people experiencing them. These processes are characterized by the adaptation of foreign goods, which entered communities through trade networks and were modified to address local concerns and to appeal to the local population. However, these processes were not uniform, as they unfolded unevenly across time, regions and social groups \parencites[p.~3-8]{hodosGlobalizationBasicsIntroduction2017}. Power dynamics are central to globalization approaches. Elites, by virtue of their social standing, had greater access to and choice in the consumption of objects, while non-elites faced constraints of limited resources and fewer options \parencites[p.~15]{pittsGlobalisationCirculationMass2014}.


An example of this approach in past assemblages is recent research on the region of Commagene in Turkey during the 1st century BCE. Earlier scholarship had interpreted the kingdom's material culture as evidence of "Hellenization", largely due to the prevalence of foreign motifs in its art and decoration. However, application of the objectscapes approach has shown that human decisions were central to the local appropriation of diverse motifs from multiple global repertoires \parencites[p.~375-379]{pittsObjectscapesManifestoInvestigating2021}. Another study by \textcite{pittsGlobalisationCirculationMass2014} has compared the spread of Roman terra \textit{sigillata} and the 17th century European consumption of Chinese porcelain within expanding trade networks. Pitts argues that "the mass consumption of long-distance ceramics as being ultimately subject to local value systems, and not simply the spread of a single global (or Roman) culture. Thus, what in strict material terms might be viewed as a process of homogenization or convergence, is often in cultural terms a re-appropriation or re-purposing of commodities to suit local needs" \parencites[p.~86]{pittsGlobalisationCirculationMass2014}.


The application of these approaches has also yielded significant insights in Sudanese archaeology, as demonstrated by recent work on New Kingdom Nubia \parencites{lemosForeignObjectsLocal2020, lemosAlternativesColonizationMarginal2021, lemosObjectMetamorphosisColonial2024}. In this context, foreign objects circulated among various Nubian communities, each using them in distinct ways. The consumption and use of these objects were shaped by social frameworks that could restrict access to certain items, particularly for non-elite groups. This highlights the need to recognize the internal diversity within Nubia. Elite and non-elite groups interacted with Egyptian-style objects differently, and their choices of adoption, rejection, and adaptation varied according to their social positions and access to resources. Consequently, objects assumed different roles depending on the local social context \parencites{lemosObjectMetamorphosisColonial2024, lemosCulturalEntanglementsExperiencing2024}. These processes created dynamic objectscapes in which Egyptian-style goods and practices were selectively incorporated, reinterpreted, and given new meaning within the Nubian cultural context, eventually integrating the local \textit{habitus}. The concept of \textit{habitus} first introduced by \textcites{bourdieuOutlineTheoryPractice1977} posits that the surrounding social and collective environment directly affects an individual's choices, becoming intrinsic to the self. Mixed families and communities of Egyptian immigrants and local Nubians would thus constitute their own \textit{habitus} and material world, shaped by choices dictated by the social collective. Therefore, objects we consider "Egyptian" in style today (hence the terms "Egyptian-style" or "Egyptianizing"), were, most likely, commonplace during the Napatan period, having been used for many generations and no longer perceived as foreign by the local population.\footnote{In this thesis, I continue to employ the term "Egyptian-style" to indicate that the amulets' motifs, for instance, were originally Egyptian.}


Funerary amulets---found across all Napatan Nubia sites and social groups---, are particularly fruitful objects for applying the objectscapes analysis. Their pervasiveness and function in funerary contexts, potentially linked to religious beliefs and cultural practices, make them ideal candidates for examining the localized interpretation and adaption of Egyptian-style designs, either imported or locally made. Give that the Napatan period encompasses burials of different social groups, multiple phases, and distinct regions within Sudanese territory---some closer to Egypt, others nearer the Napatan political center, and still others further south---, a multi-scale and material-based approach is essential for understanding amulets and uncovering objectscapes in each of these contexts. Therefore, we can investigate the origin and influence of amulet motifs, examine how they were modified to suit local preferences, and trace how iconography and use changed over time and across sites. Ultimately, the objectscapes framework has the potential to illuminate the complex interplay between Egyptian influence and local adaption, while also revealing the heterogeneity within Napatan Nubia itself, as reflected in the material choices and iconographic expressions embodied in these objects.


\section{Napatan amulets through objectscapes}


This thesis aims addresses four main research questions through the lens of the objectscapes framework:


\begin{enumerate}
	\item Can we identify patterns in the spread of  amulets in different social contexts (royal, elite and non-elite)?
	\item How did Egyptian (in origin) and local cultural aspects materialize in amulets through processes of adoption, rejection and adaptation?
	\item Do amulet designs vary across regions in Nubia? That is, are amulets from sites closer to the political center of Napata, or nearer Egypt, distinct from those further south?
	\item Can we identify longer-term evolution in amulet design throughout the Napatan period?
\end{enumerate}


In order to answer these questions, this study analyzes a dataset of funerary amulets using a combination of quantitative and qualitative methods. Quantitatively, I examine the frequency of specific iconographic motifs, material composition (e.g. faience, stones, metals), and their spatial and chronological distribution across sites and specific periods. Qualitatively, I consider the contextual associations of amulets within burial assemblages, as well as examining specific amulets and their motifs. This dual approach allows for a more nuanced understanding of how amulets were integrated into Nubian funerary customs, revealing both the impact of long-distance connectivity and the expression of local identities. By examining the circulation, transformation, and reception of these objects, this thesis contributes to a deeper understanding of cultural interaction in the Nile Valley during the Napatan period.


The analytical process of this thesis proceeds in two stages. First, Chapters 4 and 5 offer a detailed examination of the amulet assemblages from royal, elite and non-elite cemeteries through their composition, distribution, and contextual associations across these distinct social contexts. The chapters are primarily descriptive in nature, providing the empirical foundation for the subsequent analysis. Building upon this foundation, Chapter 6 incorporates the full amulet dataset, comparing amulet use across different social, regional and chronological contexts. Finally, Chapter 7 directly addresses the four research questions outlined above through the application of the objectscapes framework, examining social patterns, and regional and chronological variability. In doing so, it discusses their broader implications for understanding cultural interactions between Nubia and Egypt during the Napatan period.


Beyond its specific contributions to Sudanese archaeology, this thesis offers a case study for the broader application of the objectscapes approach in contexts of cultural contact and exchange. By  demonstrating how material culture simultaneously reflects and shapes social processes, this research highlights the importance of attending to both local agency and external influence when interpreting the archaeological record. Ultimately, the objectscapes framework enables a more dynamic and relational understanding of cultural interaction.


\section{The dataset}


This thesis is based on the analysis of a dataset built on a relational database and uses the programming language SQL. In total, there are 2,559 burials, 49,981 artifacts, and 5,480 amulets distributed across nine Nubian sites from the Napatan period. These sites span the entire territory, from Lower Nubia to beyond the Fifth Cataract in Upper Nubia.


The dataset integrates diverse sources, ensuring a comprehensive foundation for analysis. The data was taken from excavation reports published from the beginning of the 20th century to the 1980s, as well as modern analysis \parencites{griffithOxfordExcavationsNubia1923, dunhamElKurru1950, dunhamNuri1955, dunhamWestSouthCemeteries1963, geusCimetièreXXVèmeDynastie1975, vilaNécropoleMissiminiaSépultures1980, säve-söderberghScandinavianJointExpedition1989, säve-söderberghScandinavianJointExpedition1989a, williamsExcavationsAbuSimbel1990, lohwasserAspekteNapatanischenGesellschaft2012}. Sanam, in particular, never had its full reports published, but \textcite{lohwasserAspekteNapatanischenGesellschaft2012}, in her analysis of the site, transcribed Griffith’s complete description of all burials.

The database does not include only amulets but also artifacts, as they were used to analyze trends across phases as well as distinctions between social status. Firstly, the data will be analyzed based on the social groups represented in these burials—royal, elite and non-elite (Chapters 4 and 5)—, followed by a comparative analysis between the different social groups focusing specifically on amulets (Chapter 6).


Additionally, efforts were made to locate the amulets in museums and include their register numbers in the database. For Sanam, many amulets were already located by Lohwasser for her publication of the site. Furthermore, numerous amulets were catalogued in the publication by \textcites{kormyshevaGodsDivineSymbols2006}, which focused on the Sudan National Museum collection. Some amulets from Mirgissa were found thanks to the Global Egyptian Museum’s online catalogue. \footnote{\href{http://www.globalegyptianmuseum.org}{www.globalegyptianmuseum.org}}

The abbreviations for the museums used in the database are as follows:

\begin{itemize}
	\item Ash.: Ashmolean Museum of Art and Archaeology
	\item ÄMP: Ägyptischen Museums und Papyrussammlung, Berlin
	\item Brigh.: Brighton \& Hove Museums
	\item Brüssels: Royal Museums of Art and History, Brussels
	\item Cairo: The Egyptian Museum in Cairo
	\item Fitz.: The Fitzwilliam Museum
	\item Leiden: Rijksmuseum van Oudheden (National Museum of Antiquities)
	\item Manchester: Manchester Museum
	\item MFA: Museum of Fine Arts, Boston
	\item Penn Museum: The University of Pennsylvania Museum of Archaeology and Anthropology
	\item Pitt.: Pitt Rivers Museum
	\item SNM: Sudan National Museum
\end{itemize}


Finally, the complete database will be made available online for consultation, providing researchers and future scholars with an accessible resource for further study.


\section{Possibilities and limitations of the dataset}


\subsection{General considerations}


Unfortunately, almost all of the cemeteries examined in this thesis have suffered disturbance due to looting, most likely occurring multiple times since antiquity. This is not only the case for older excavations, which already lacked modern archaeological methods, but also more recent excavations, which have been affected by previous disturbances. Royal tombs were often the primary targets of looters due to their imposing architecture and valuable contents. Furthermore, tomb superstructures have frequently been eroded or destroyed by later constructions, complicating the identification of tomb contexts, while tomb reuse in the subsequent Meroitic period is also attested across sites, further concealing the original Napatan contexts. 


The extent of disturbance is unequivocal across all sites, yet the available reports provide inconsistent and often incomplete information. At the royal cemeteries of El-Kurru and Nuri, Dunham explicitly status that none were intact \parencites[p.~12]{dunhamElKurru1950}[p.~4]{dunhamNuri1955}. At Meroe West and South, the reports only occasionally remark on a tomb's intact status, sometimes noting that the blocking to the burial chamber remained undisturbed. Griffith's brief publication on Sanam states that all tombs were plundered except for the "poorest ones", although even in Lohwasser's comprehensive documentation of his field report, there is no systematic designation of which tombs were disturbed and which were not \parencites[p.~75]{griffithOxfordExcavationsNubia1923}[p.~423-509]{lohwasserAspekteNapatanischenGesellschaft2012}. At Qustul, Williams only comments on the disturbance of the body itself, without addressing whether other parts of the burial were also affected. For Debeira 176, Säve-Soderbergh note in the main text that "the site was much disturbed by plundering" \parencite[p.~200]{säve-söderberghScandinavianJointExpedition1989}, while the accompanying table records whether the superstructure, burial place, and/or human remains were disturbed. Although some entries were left blank, none were explicitly recorded as undisturbed \parencite[p.~36-38]{säve-söderberghScandinavianJointExpedition1989a}. At Mirgissan, Geus observes that, despite an apparent absence of widespread looting, few tombs were found intact, because most contained very few objects or were completely empty. He also notes that the site suffered extensive erosion, which may contributed for the poor preservation \parencite[p.~479]{geusCimetièreXXVèmeDynastie1975}. For Missiminia, the tombs suffered the same fate, with Vila noting that erosion destroyed many aspects of the burials and that almost all tombs showed evidence of plundering \parencite[p.~11, 26-29, 32]{vilaNécropoleMissiminiaSépultures1980}. Given this inconsistent and incomplete reporting, it is impossible to determine the level of disturbance on an individual basis. No standardized criteria was applied across excavations, and there is no way to assess whether one tomb was more or less affected than another. Therefore, in the analysis, some level of disturbance has to be accounted for. While this limits the types of conclusions that can be drawn from quantitative comparisons, the analysis of patterns across social, chronological, and regional dimensions remains a viable and productive approach.


A further challenge concerns the identification and determination of sex and gender from human remains\footnote{I understand sex as encompassing biological and physiological differences, although osteological identification of sex in ancient human remains is not straightforward, encompassing a spectrum from extremely female to extremely male with many variations in between. Gender, by contrast, concerns social identity. See \textcite[p.~385-398]{whiteHumanBoneManual2005}}. The disturbances that affected the cemeteries have also resulted in damage and even the absence of many human remains, making osteological analysis impossible in most cases. Moreover, osteological analysis was not yet a formalized field of study at the time of the older excavations, and none of the excavation reports specify the methodology by which the tomb occupants were identified. Consequently, the sex determinations provided in the reports must be treated with considerable caution. 


The reliability of sex and gender attribution varies considerably depending on the type of burial. In the case of royal burials (and one elite tomb), gender seems to have been assigned by excavators mostly based on non-osteological criteria, such as inscriptions, the types of shabtis, tomb architecture, and the tomb's spatial relation to that of a king (in the case of queens and minor wives).\footnote{Although Reisner drew some conclusions based on the human remains, subsequent analysis did not correlate the sex of the body with the presumed gender \parencite{beckDemographicDataHuman1999}.} For the remaining burials, sex determination are scarcely included, and where they do appear (Sanam, Qustul, Debeira 176, and the Meroe cemeteries), they lack any methodological specifications. Given these inconsistencies, the sex data recorded in the database reflect the available and published attributions by the excavators but cannot be considered reliable for analytical purposes. Consequently, it will not be used as a variable in this thesis, except in the case of royal burials, where attributions of gender and royal title were based on material evidence.


Another limitation pertains to the documentation of amulets. Some descriptions were very limited, without providing additional details apart from the main motif. This is further complicated by the absence of photographs and drawings, particularly in non-elite contexts, which hinders not only the identification of complementary motifs and previously unidentified designs. In order to resolve this, attempts were made to locate these amulets in museums or other institutions where they could be currently housed. Unfortunately, many cases lacked sufficient information to complete the search. Consequently, in many cases, I had to rely solely on the description provided in the reports.


\subsection{Chronological division}


Another methodological limitation is the chronology of the period, specifically precise dating of individual elite and non-elite tombs. For royal tombs, a chronological division of the Napatan period was established by Reisner, who separated the period into generations of kings, along with the burials of their associated queens and minor wives \parencites{dunhamElKurru1950, dunhamNuri1955}. \footnote{In the site reports, the female burials are classified as either "queen or (unnamed) woman", but for this thesis I chose to use the term "minor wife", which better captures the status of these women as secondary consorts to the king, with the queen being their main wife.} This division was used for the elite and non-elite graves at Meroe West and South. However, since accurate dating of these tombs is more complex, Reisner assigned them broader dates, spanning many generations.


For this thesis, the Napatan period was divided into five phases: pre-25th Dynasty, 25th Dynasty, Early Napatan, Middle Napatan, and Late Napatan (Table \ref{tab:phases-divison}). . This broader division was based on the historical and political contexts of each phase, which may have impacted the use of amulets.\footnote{The historical background of the Napatan period will be further explored in Chapter 3.} Royal tombs were assigned their phases based on this division, as well as elite and non-elite graves.


The pre-25th Dynasty phase encompasses the time between the end of the Egyptian rule of Lower Nubia and the start of the Napatan dynasty―considered to have been established by Alara \parencites[p.~314-318]{törökTwoWorldsFrontier2009}. The 25th Dynasty comprehends the kings that expanded the Kushite rule into Upper Egypt and reigned as kings in both territories. After their withdrawal from Egypt by the Assyrians, four rulers were included in the Early Napatan phase. This phase is characterized by continued trade and religious ties between Thebes and Napata. Despite these relations, a military campaign was carried out by the Assyrian ruler in Egypt targeting Lower Nubia. Kings from the Middle Napatan period did not leave textual evidence, however, trades with Persia were documented in foreign texts. In contrast, the Late Napatan phase included kings known mainly through stelae, temple inscriptions and their pyramids. This broader division was based on the historical and political contexts of each phase, which may have impacted the use of amulets. Royal tombs were assigned their phases based on this division, as well as elite and non-elite graves.


{\small
	\begin{xltabular}{\textwidth}{X X X}
		\caption{Chronological division of the Napatan period and its associated owner and tomb.}
		\label{tab:phases-divison} \\
		\toprule
		\textbf{Owner (if known)} & \textbf{Burial place} \\
		\midrule
		\endfirsthead
		\multicolumn{2}{c}{\textit{Continued from previous page}} \\
		\toprule
		\textbf{Owner (if known)} & \textbf{Burial place} \\
		\midrule
		\endhead
		\bottomrule
		\endlastfoot
		% --- DATA ---
		\multicolumn{2}{c}{\textbf{Pre-25th Dynasty}} \\
		\midrule
		King & Ku. Tum. 1 \\
		King & Ku. Tum. 2 \\
		King or queen & Ku. Tum. 4 \\
		King or queen & Ku. Tum. 5 \\
		King & Ku. Tum. 6 \\
		King or queen & Ku. 19 \\
		King & Ku. 14 \\
		King or queen & Ku. 13 \\
		King & Ku. 11 \\
		Queen (?) & Ku. 10 \\
		King Alara & Ku. 9 \\
		Queen Kasaqa & Ku. 23 \\
		Minor wife (?) & Ku. 21 \\
		\midrule
		\multicolumn{2}{c}{\textbf{25th Dynasty}} \\
		\midrule
		King Kashta & Ku. 8 \\
		Queen Pebatma (?) & Ku. 7 \\
		Minor wife (?) & Ku. 20 \\
		King Piankhy & Ku. 17 \\
		Queen (Peksater) (?) & Ku. 22 \\
		Queen (?) & Ku. 55 \\
		Minor wife (Peksater) (?) & Ku. 54 \\
		Queen Neferukakashta & Ku. 52 \\
		Queen (?) & Ku. 51 \\
		Queen Tabiry & Ku. 53 \\
		Queen Khensa & Ku. 4 \\
		King Shabataqo & Ku. 18 \\
		Queen (?) & Ku. 72 \\
		Minor wife (?) & Beg. S 155 \\
		King Shabaqo & Ku. 15 \\
		Queen (?) & Ku. 71 \\
		Queen (?) & Ku. 62 \\
		Minor wife & Beg. W 609 \\
		King Taharqa & Nu. 1 \\
		Queen Naparaye & Ku. 3 \\
		Queen (Abar) (?) & Nu. 35 \\
		Queen Atakhebasken & Nu. 36 \\
		King Tanwetamani & Ku. 16 \\
		Queen Qalhata & Ku. 5 \\
		Queen Arty (?) & Ku. 6 \\
		Queen Malaqaye & Nu. 59 \\
		Minor wife (?) & Nu. 74 \\
		Minor wife (?) & Nu. 77 \\
		Minor wife (?) & Nu. 80 \\
		\midrule
		\multicolumn{2}{c}{\textbf{Early Napatan}} \\
		\midrule
		King Atlanersa & Nu. 20 \\
		Queen (?) & Ku. 61 \\
		Queen Yeturow & Nu. 53 \\
		Queen (?) & Nu. 60 \\
		Minor wife (?) & Nu. 75 \\
		King Senkamanisken & Nu. 3 \\
		Queen (Maletaral I) (?) & Nu. 41 \\
		Minor wife & Nu. 71 \\
		Minor wife (?) & Nu. 73 \\
		Minor wife (?) & Nu. 78 \\
		Minor wife (?) & Nu. 81 \\
		Minor wife & Nu. 82 \\
		King Anlamani & Nu. 6 \\
		Queen (Takahatenamun) (?) & Nu. 21 \\
		Queen (Amanimalel) (?) & Nu. 22 \\
		Queen Masalaye & Nu. 23 \\
		Minor wife (?) & Nu. 72 \\
		Minor wife (?) & Nu. 76 \\
		Minor wife (?) & Nu. 79 \\
		King Aspelta & Nu. 8 \\
		Queen Nasalsa & Nu. 24 \\
		Queen Madiqen & Nu. 27 \\
		Queen Maqemale & Nu. 40 \\
		Queen Asata & Nu. 42 \\
		Queen Artaha & Nu. 58 \\
		Queen/princess Mernua & Beg. S 85 \\
		\midrule
		\multicolumn{2}{c}{\textbf{Middle Napatan}} \\
		\midrule
		King Aramatelqo & Nu. 9 \\
		Queen Henuttakhebit & Nu. 28 \\
		Queen Akheqa & Nu. 38 \\
		Queen Maletasen & Nu. 39 \\
		Minor wife (?) & Nu. 54 \\
		Queen Atmataka & Nu. 55 \\
		Queen Piankh-her (?) & Nu. 57 \\
		King Malonaqen & Nu. 5 \\
		Queen Amanitakaye & Nu. 26 \\
		Queen Tagtal (?) & Nu. 45 \\
		King Analmakheye & Nu. 18 \\
		King Amani-nataki-lebte & Nu. 10 \\
		Queen Maletaral & Nu. 25 \\
		King Karkamani & Nu. 7 \\
		Queen (?) & Nu. 30 \\
		King Amaniastabarqo & Nu. 2 \\
		Minor wife (?) & Nu. 47 \\
		Minor wife (?) & Nu. 50 \\
		King Sikhespiqo & Nu. 4 \\
		Queen Piankh-qe-qa & Nu. 29 \\
		Minor wife (?) & Nu. 52 \\
		King Nasakhma & Nu. 19 \\
		Minor wife (?) & Nu. 46 \\
		Minor wife (?) & Nu. 49 \\
		King Malewi-amani & Nu. 11 \\
		Queen Sakachaye & Nu. 31 \\
		Queen Akhrasan & Nu. 32 \\
		King Talakhamani & Nu. 16 \\
		King Irike-Amanote & Nu. 12 \\
		Queen (?) & Nu. 33 \\
		\midrule
		\multicolumn{2}{c}{\textbf{Late Napatan}} \\
		\midrule
		King Baskakeren & Nu. 17 \\
		King Harsijotef & Nu. 13 \\
		Queen Henutirdis & Nu. 34 \\
		Queen Batahaliye & Nu. 44 \\
		Queen Atasamale & Nu. 61 \\
		King (unused) & Ku. 1 \\
		Queen & Ku. 2 \\
		Minor wife (?) & Nu. 43 \\
		Minor wife (?) & Nu. 48 \\
		Minor wife (?) & Nu. 51 \\
		King Akhraten & Nu. 14 \\
		Queen (?) & Nu. 37 \\
		King Nastasen & Nu. 15 \\
		Queen (Sakhmakh or Pelkha) (?) & Nu. 56 \\
	\end{xltabular}
}


\subsubsection{Categorization of object types}


Given the large number of objects types found in the burials, for a better visualization and analysis, some were divided into broader categories of objects. These categories will be used throughout the thesis. Below is a list of all of the categories used in the database. Each category will be described in more detail if necessary.


\begin{itemize}
	\item Adornments: objects to be used on the body such as bracelets, necklaces, collars, girdles, rings, earrings, etc. This category includes the whole object, with beads being a separate category. Beads were found in large numbers loose on the burial and, therefore, it would askew the data to consider each individual bead as a separate adornment. 
	\item Animal parts: this category comprises the bones of various animals found within the burial, including ox, calf, camel, dog and horse. It also encompasses other animal remains, such as horns, teeth, feathers, eggshells—particularly from ostriches and geese—, as well as shells like cowrie and snail shells. The latter were only included if explicitly described as unpierced; pierced shells were instead classified under beads or adornments, with the shell type specified as the material.
	\item Baskets: includes baskets primarily made of organic material. They were only found in Meroe West.
	\item Bead nets: encompasses objects identified in the reports as either bead or mummy nets.
	\item Beads: as noted earlier, this category only includes loose beads, excluding those incorporated into adornments or bead nets.
	\item Burial containers and trappings: this category contains complete containers (e.g. coffins and beds), as well as detached components (e.g. coffin lids and bed legs). Associated fittings and decorative inlays (e.g. coffin fittings and eye/eyebrow inlays) are likewise included.
	\item Canopics: all objects designated by Reisner/Dunham as canopic-related are included here, encompassing jars (standard and miniature), lids, figures of the Four Sons of Horus, canopic box and figure fittings.
	\item Ceramic vessels: this classification included all kinds of vessels (e.g. bowls, jars, pots, amphorae, etc) made of pottery.
	\item Clothing: includes objects described in the reports as sandals, shawls, wraps and garments.
	\item Toiletries: objects related to the cosmetic and personal care sphere, including kohl, tools to apply and store kohl, cosmetic spoons, mirrors and other toilet instruments, some not specified.
	\item Cylinder sheaths: a type of object found only in Nuri. They will be further explored in the thesis.
	\item Decorative elements: this category includes items explicitly described in the reports as tomb decorations, as well as those lacking a clear association with a specific object but presumed to be decorative in nature. Therefore, this decorative items could be from either the tomb or objects.
	\item Fragments of materials: this group consists of material pieces and fragments mentioned in the reports that lack a recognizable form.
	\item Furniture and trappings: this category encompasses both complete objects (e.g. boxes, tables, canopy poles, folding stools), as well as associated parts (e.g. decorations, legs, locks).
	\item Games and toys: includes objects associated with recreation, such as game pieces, dice and toys.
	\item Headrests.
	\item Heart scarabs.
	\item Incense burners.
	\item Jar stoppers: this category comprises objects made from clay, pottery, or plaster, designed to seal vessels (often containing goods). Some examples bear seal or thumb impressions.
	\item Lamps: this group includes objects described as either lamps or saucer lamps.
	\item Miscellaneous: comprises objects that defy standard categorization due to ambiguous use or incomplete identification in the reports. Notable examples include bells and a singular steatite replica of a Nile tortoise shell.
	\item Mummy trappings: this category includes items related to mummies, such as finger caps, masks and collars.
	\item Non-ceramic vessels: includes any kind of vessel not made from pottery (e.g. faience, travertine, silver, glass, gold, etc).
	\item Offering tables.
	\item Offering vessels: this category comprises all vessels identified in the reports as being used for offerings, regardless of material composition.
	\item Osiris whip: a unique type of object with only one example in Nuri. It is a model of a whip, associated with the god Osiris, made of gold.
	\item Ostraca.
	\item Palettes.
	\item Pebbles.
	\item Scarabs.
	\item Sealing: this category includes items associated with sealing, such as seals (which were not described as scarabs) and seal impressions.
	\item Shabtis.
	\item Shrines: this group comprises any type of shrine (e.g standard and miniature).
	\item Staffs and trappings: this category comprises both complete staffs, as well as parts of it (e.g. caps and ends).
	\item Statuettes: includes small sized statues of animals, deities, humans, among others. They differ from amulets as they did not include suspension rings and were normally bigger.
	\item Stelae.
	\item Tools: this broad category encompasses functional objects for crafting, processing, or professional or daily tasks.
	\item Weapons: this category includes objects associated with combat and hunting activities. Included items range from complete weapons to functional components, such as arrows, bows, spears and daggers.
\end{itemize}


\section{Distinguishing social status}


The most important research question of this thesis is the analysis of amulets from burials of various social groups. To accurately differentiate between these burials, specific criteria have been utilized for this study. This is particularly important for the differentiation of elite and non-elite burials from the same sites. Royal cemeteries, in particular, offer clear evidence of this differentiation due to the inscribed material that identifies kings and queen buried there. Even when inscriptions are not present, the location of the tombs was used to estimate the possible owner. Given these established royal burial sites, it is improbable that non-royal individuals would have been permitted to be interred within them.


In Egyptological literature, a few studies have aimed to establish a framework through several factors in order to distinguish the social status of individuals based on their burial sites and artifacts. Overall, they include the following criteria: level of human and material resource in tomb construction, quantity of grave goods, presence of burial-specific objects―such as body containers, shabtis and canopics―, presence of inscriptions, and objects crafted from "prestigious" materials \parencites{smithIntactTombsSeventeenth1992, richardsSocietyDeathAncient2005, wadaProvincialSocietyCemetery2007, podvinLéquipementTombesIntactes2017, lemosMaterialCultureSocial2017}.


By applying these criteria to the database, primarily to differentiate elite and non-elite burials located in the same areas, we can better understand the social dynamics reflected in their burials practices. Table \ref{tab:social-group-division} correlates the different elements: superstructures, substructures, and the median number of burial-specific objects, inscribed objects and those made of prestigious materials. For this analysis, I am employing the median instead of average. The average is misleading here because many tombs were completely empty, which pulls the data downward. The median, by contrast, identifies the middle value in the dataset and is therefore unaffected by the extremes. In order to get a true picture of how these objects are spread across contexts, the median is a better metric.


Based on the database, five types of superstructure were identified: double pyramids\footnote{As we will see in Chapter 4, these pyramids were build in two stages, with the second structure erected after the death of the owner, underscoring their continuous prestige.}, pyramids, tumuli, unidentified (either non-existent or destroyed), and no superstructure at all. Among the substructures there are: chambers, cave tombs, and pits. The constructions which required higher access and consumption of human and material resources, are both types of pyramids and chambers.


Therefore, it would be expected that these tombs would yield the most median number per tomb in all criteria of artifacts. However, as shown in Table \ref{tab:social-group-division}, the relationship between tomb structures and the artifacts is more complex.


Pyramids with chambers take the lead in only the highest median number of total objects and inscribed items. However, they did not lead in categories such as burial-specific objects or those made of prestigious materials. \footnote{Materials considered prestigious are gold, silver, lapis lazuli, obsidian and ivory. Gold and ivory were commodities of Nubia, while lapis lazuli and obsidian had to be imported from Afghanistan and Yemen, respectively. Silver could be found in gold deposits or in the Egyptian Eastern Desert.} Another tomb structure with high consumption of resources are double pyramids, which had high medians across all criteria, confirming their correlation.


However, tombs with a pit but without a superstructure (most likely eroded), contained the same median of artifacts and burial-specific objects as pyramids with pit. In fact, these tombs also contained more grave goods than chambers without a superstructure.
had the most median of objects, while with pits the numbers were comparable to no superstructure with chambers and pits. In fact, no superstructure with pits contained a bigger median of objects, inscribed objects and objects made of prestigious materials than those with chambers.


Therefore, it is not totally straightforward to assign the status of elite or non-elite based on the correlation of the criteria above. Especially considering that almost all of the tombs were plundered and pyramids were the primary target. Consequently, a complete analysis of the objects is not possible. Therefore, for this thesis, I have adopted a resource expenditure approach, in which tombs with the higher need for human and material resources were considered elite and the ones with less resources needed as non-elite. This distinction will be further elaborated in Chapter 5, where these burials will be analyzed in more detail.


{\small
	\begin{xltabular}{\textwidth}{X X X X X}
		\caption{Correlation between tomb super- and substructures with the median quantity of objects, burial-specific objects, inscribed objects and objects made of prestigious materials.}
		\label{tab:social-group-division} \\
		\toprule
		\textbf{Substructures} & \textbf{Median of objects} & \textbf{Median of burial-specific objects} & \textbf{Median of inscribed objects} & \textbf{Median of prestigious material objects} \\
		\midrule
		\endfirsthead
		\multicolumn{5}{c}{\textit{Continued from previous page}} \\
		\toprule
		\textbf{Substructures} & \textbf{Median of objects} & \textbf{Median of burial-specific objects} & \textbf{Median of inscribed objects} & \textbf{Median of prestigious material objects} \\
		\midrule
		\endhead
		\bottomrule
		\endlastfoot
		\multicolumn{5}{c}{\textbf{Pyramid}} \\
		\midrule
		Chambers & 59 & 2 & 37.5 & 2 \\
		Pit & 7 & 1 & 2 & 3 \\
		\midrule
		\multicolumn{5}{c}{\textbf{Double Pyramid}} \\
		\midrule
		Chambers & 45 & 6 & 16 & 8 \\
		\midrule
		\multicolumn{5}{c}{\textbf{Unknown}} \\
		\midrule
		Cave tomb & 2 & 1 & 1 & 2 \\
		Chambers & 27.5 & 2 & 3 & 2.5 \\
		Pit & 5 & 1 & 2 & 3 \\
		\midrule
		\multicolumn{5}{c}{\textbf{No superstructure}} \\
		\midrule
		Chambers & 3.5 & 2 & 0 & 1.5 \\
		Pit & 7 & 1 & 1 & 2 \\
		\midrule
		\multicolumn{5}{c}{\textbf{Tumulus}} \\
		\midrule
		Pit & 4 & 0 & 1 & 1 \\
	\end{xltabular}
}